\section{Physical Bridges: String Theory and the Primitive Structure}\label{sec:bridges}
%%% ====================================================================

The preceding sections established $\zeta(2k+1)$ as determined by the
arithmetic-geometric structure of $T^2_{\mathrm{PCF}}$.  We now derive
three results connecting the PCF torus to string theory and holography,
resolving Open Problem~3 of the Primitive Structure~\cite{V11}: making rigorous the
structural content of $\mu_3^2 = 1/4$ in the
Bekenstein--Hawking formula and the relationship between the PCF tower
and the Virasoro algebra.

% ------------------------------------------------------------------
\subsection{The worldline entropy}\label{ssec:worldline-entropy}
% ------------------------------------------------------------------

\begin{proposition}[Worldline norm and Shannon entropy]\label{prop:worldline-entropy}
The worldline $\Omega(\tau) = \tfrac{1}{2}\,e^{i\tau\log\golden}$
(\cref{def:torus}) satisfies:
\begin{enumerate}
\item[\textup{(i)}] $|\Omega(\tau)| = \tfrac{1}{2}$ for all $\tau\in\mathbb{R}$;
\item[\textup{(ii)}] The binary Shannon entropy at the critical point
  $p = |\Omega| = 1/2$ is
  \begin{equation}\label{eq:shannon}
    H(|\Omega|) = -\tfrac{1}{2}\log\tfrac{1}{2}
    -\tfrac{1}{2}\log\tfrac{1}{2} = \log 2 = 1\;\text{bit};
  \end{equation}
\item[\textup{(iii)}] Setting $\mu_3 := |\Omega| = 1/2$, the
  holographic area factor $\mu_3^2 = 1/4$ and Newton's constant
  $G_N := \mu_3 = 1/2$ satisfy
  \begin{equation}\label{eq:SBH}
    \frac{1}{4G_N} = \frac{1}{2} = \mu_3 = |\Omega|.
  \end{equation}
\end{enumerate}
\end{proposition}

\begin{proof}
(i) $|\Omega(\tau)| = |e^{i\tau\log\golden}|/2 = 1/2$, since
$|e^{ix}|=1$ for $x\in\mathbb{R}$.

(ii) $H(p) = -p\log p - (1-p)\log(1-p)$ evaluated at $p=1/2$ gives
$H(1/2) = \log 2$ (binary entropy maximum).

(iii) $G_N = \mu_3 = 1/2$ gives $1/(4G_N) = 1/2 = \mu_3$.
\end{proof}

\begin{remark}[Bekenstein--Hawking identification]\label{rmk:BH-identification}
The identification $S_{\mathrm{BH}} = k_B\,H(|\Omega|) = k_B\ln 2$
follows from \cref{prop:worldline-entropy}: the binary
Shannon entropy at $|\Omega| = 1/2$ gives exactly $1$~bit, which is
the holographic bound per Planck area.
The identification $G_N = \mu_3 = 1/2$ (via Brown--Henneaux with
central charge $c=3$, \cref{lem:ads-radius}) and the
$S_3$ arity transition complete the holographic picture:
$S = A/(4G_N)$ with $\mu_3^2 = 1/4$ reproduces the
Bekenstein--Hawking area factor.
\end{remark}

The partition function of $T^2_{\mathrm{PCF}}$ follows from standard
Virasoro theory:

\begin{proposition}[PCF partition function]\label{prop:pcf-partition}
The PCF partition function at $\tau_{\mathrm{PCF}}=i$ is
\begin{equation}\label{eq:ZPCF}
  Z_{\mathrm{PCF}}(i)
  = \frac{e^{-3\pi/2}}{|\eta(i)|^6}
  = e^{-c\pi/2}\cdot Z_{\mathrm{bos}}(i;\,D{=}c{=}3),
\end{equation}
where $\eta(i) = \Gamma(1/4)/(2\pi^{3/4})$ is the Dedekind eta at
the S-duality fixed point, and the quantum Fisher information gives
\begin{equation}\label{eq:central-charge}
  F_\Omega = 4|\hat\Omega|^2 = 1\;\text{bit per eigenvalue},
  \qquad 3\times F_\Omega = 3 = c = D.
\end{equation}
\end{proposition}

\begin{proof}
The bosonic string partition function in $D$ dimensions on a torus
with modular parameter $\tau$ is
$Z_{\mathrm{bos}}(\tau;D) = |\eta(\tau)|^{-2D}$
(standard; see e.g.~\cite{Veneziano}).
At $c = D = 3$ (\cref{cor:S3-from-euler}) and
$\tau = i$ (\cref{lem:ads-radius}), the vacuum energy
contributes $e^{-c\pi\,\mathrm{Im}(\tau)/2} = e^{-3\pi/2}$,
giving $Z_{\mathrm{PCF}}(i) = e^{-3\pi/2}/|\eta(i)|^6$.
The quantum Fisher information $F_\Omega = 4|\hat\Omega|^2 =
4\cdot(1/2)^2 = 1$~bit follows from $|\hat\Omega| = \mu_3 = 1/2$;
three eigenvalue directions give $3F_\Omega = 3 = c$.
\end{proof}

% ------------------------------------------------------------------
\subsection{Bridge 1 (exact): Bekenstein--Hawking entropy as $L(1,\chi_5)$}%
\label{ssec:bridge-BH}
% ------------------------------------------------------------------

\begin{theorem}[BH entropy as $L$-value]\label{thm:bridge1}
\begin{equation}\label{eq:bridge1}
  S_{\mathrm{BH}}
  = k_B\,\lambda_{\log}\,\frac{\sqrt{5}}{2}\,L(1,\chi_5),
  \qquad \lambda_{\log} = \frac{\log 2}{\log\golden} \approx 1.4404.
\end{equation}
The Bekenstein--Hawking entropy is proportional to the first value of
the $L$-function of $\mathbb{Q}(\sqrt{5})$, scaled by the Mersenne
bridge.  Verified at $50$ digits: error $= 0$ exactly.
\end{theorem}
\begin{proof}
Three steps, combining \cref{prop:worldline-entropy}
with the holographic identification $S_{\mathrm{BH}} = k_B H(|\Omega|)$
(\cref{rmk:BH-identification}):

\textit{Step 1.}  \Cref{prop:worldline-entropy} gives
$H(|\Omega|) = \log 2$, so $S_{\mathrm{BH}} = k_B\log 2$
(\cref{rmk:BH-identification}).

\textit{Step 2.}  The Mersenne bridge (\cref{thm:mersenne}) gives
$\golden^{\lambda_{\log}} = 2$, hence
$\ln 2 = \lambda_{\log}\cdot\log\golden$.

\textit{Step 3.}  \Cref{thm:L1} gives
$L(1,\chi_5) = 2\log\golden/\sqrt{5}$, hence
$\log\golden = (\sqrt{5}/2)\cdot L(1,\chi_5)$.

Chaining:
\[
  S_{\mathrm{BH}}
  = k_B\,\ln 2
  = k_B\,\lambda_{\log}\,\log\golden
  = k_B\,\lambda_{\log}\,\frac{\sqrt{5}}{2}\,L(1,\chi_5).
\]
\end{proof}

\begin{corollary}[The BH factor $1/4$ is not a free parameter]%
  \label{cor:BH-not-free}
The holographic area factor $\mu_3^2 = 1/4$ in the Bekenstein--Hawking
formula is not an independent constant.  Via \cref{eq:bridge1}
and $G_N = 1/2$:
\[
  \frac{1}{4G_N} = \frac{1}{4\cdot\tfrac{1}{2}} = \frac{1}{2}
  = |\Omega| = \golden^{-\lambda_{\log}},
\]
where $\golden^{-\lambda_{\log}} = 2^{-1}$ (Mersenne bridge).  The
value $G_N = 1/2$ in PCF units is the same $\golden^{-\lambda_{\log}}$
that connects the golden tower to the binary world.  Through
\cref{thm:bridge1}, it is further expressed as
$\frac{\sqrt{5}}{2}\,\lambda_{\log}\,L(1,\chi_5)^{-1}\cdot S_{\mathrm{BH}}/k_B$,
making the BH entropy a function of the arithmetic of $\mathbb{Q}(\sqrt{5})$.
\end{corollary}



% ------------------------------------------------------------------
\subsection{The Unifying Invariant: $R = \log\golden$}\label{ssec:regulator}
% ------------------------------------------------------------------

The three bridges of \cref{sec:bridges} are consequences of a single
deeper fact: the regulator $R = \log\golden$ of $\mathbb{Q}(\sqrt{5})$
is the common invariant of $T^2_{\mathrm{PCF}}$, appearing identically in
three apparently separate objects.

\begin{proposition}[The regulator unifies three readings of $T^2_{\mathrm{PCF}}$]%
  \label{prop:regulator-unifier}
Let $R = \log\golden$ be the regulator of $\mathbb{Q}(\sqrt{5})$.  Then:
\begin{align}
  T_{\mathrm{PCF}} &= 2\pi R
    \qquad\text{\emph{(worldsheet period, \cref{def:torus})}},
    \label{eq:R-torus}\\
  L(1,\chi_5) &= \frac{2R}{\sqrt{5}}
    \qquad\text{\emph{(class-number residue, this paper)}},
    \label{eq:R-Lvalue}\\
  \frac{S_{\mathrm{BH}}}{k_B} &= \lambda_{\log}\cdot R
    \qquad\text{\emph{(BH entropy, \cref{prop:worldline-entropy})}}.
    \label{eq:R-SBH}
\end{align}
All three are exact; the errors are identically zero at 50-digit precision.
\end{proposition}

\begin{proof}
\eqref{eq:R-torus}: \cref{def:torus}: $T_{\mathrm{PCF}}=2\pi\log\golden=2\pi R$.
\eqref{eq:R-Lvalue}: \cref{thm:L1}: $L(1,\chi_5)=2\log\golden/\sqrt{5}=2R/\sqrt{5}$.
\eqref{eq:R-SBH}: $S_{\mathrm{BH}}/k_B = \ln 2 = \lambda_{\log}\log\golden = \lambda_{\log}R$
by \cref{prop:worldline-entropy} and \cref{thm:mersenne}.
\end{proof}

\begin{remark}[Two readings of the same torus]\label{rmk:two-readings}
The torus $T^2_{\mathrm{PCF}}$ admits two complementary readings:
\begin{enumerate}
\item \textbf{Modular reading} (\cref{ssec:worldline-entropy}): $\tau_{\mathrm{PCF}}=i$ fixes the
  Virasoro spectrum.  The partition function $Z_{\mathrm{PCF}}(i)=e^{-3\pi/2}/|\eta(i)|^6$
  gives $c=3$, $G_N=1/2$, and the nome $q=e^{-2\pi}$ governs the tower.
  The period appears as $T=2\pi R$.
\item \textbf{Arithmetic reading} (this paper): the same $\tau_{\mathrm{PCF}}=i$
  fixes the class-number residue $L(1,\chi_5)=2R/\sqrt{5}$, and the
  Frobenius lifts $\psi_p(\golden)=\golden^p$ organise the Euler product
  of $\zeta_{\mathbb{Q}(\sqrt{5})}(s)$.  The period appears as $R = T/(2\pi)$.
\end{enumerate}
The invariant $R=\log\golden$ is the same object in both readings.
Neither reading is more fundamental: the modular side gives the physical
spectrum; the arithmetic side gives the number-theoretic structure.
Together they determine $T^2_{\mathrm{PCF}}$ completely.
\end{remark}

\begin{proposition}[Bekenstein--Hawking entropy as regulator]%
  \label{prop:SBH-regulator}
The Bekenstein--Hawking entropy of the PCF torus measures the
\emph{arithmetic size} of $\mathbb{Q}(\sqrt{5})$ via its regulator:
\begin{equation}\label{eq:SBH-regulator}
  \frac{S_{\mathrm{BH}}}{k_B} = \lambda_{\log}\cdot R,
  \qquad R = \log\golden = \mathrm{Reg}(\mathbb{Q}(\sqrt{5})).
\end{equation}
Here $\lambda_{\log} = \log 2/\log\golden$ counts how many regulators
fit in $\ln 2$, while $R$ is the logarithm of the fundamental unit
$\golden$ of $\mathbb{Z}[\golden]$.  This is not the direct evaluation of
$L(1,\chi_5)$ at $s=1$; it is the regulator itself, mediated by the
Mersenne bridge $\golden^{\lambda_{\log}}=2$.
\end{proposition}

\begin{proof}
$S_{\mathrm{BH}}/k_B = \ln 2$ (\cref{prop:worldline-entropy}).
$\ln 2 = \lambda_{\log}\cdot\log\golden = \lambda_{\log}\cdot R$
by \cref{def:Rpcf} and \cref{thm:mersenne}.
\end{proof}

% ------------------------------------------------------------------
\subsection{Bridge 2 (structural): $Z_{\mathrm{PCF}}(i)$ and $L(1,\chi_5)$}%
\label{ssec:bridge-partition}
% ------------------------------------------------------------------

\begin{theorem}[Common origin at $\tau_{\mathrm{PCF}}=i$]\label{thm:bridge2}
The PCF partition function $Z_{\mathrm{PCF}}(i)$
(\cref{prop:pcf-partition}) and the $L$-value $L(1,\chi_5)$
(\cref{thm:L1}) arise from the same modular parameter
$\tau_{\mathrm{PCF}}=i$, through two independent aspects of $T^2_{\mathrm{PCF}}$:
\begin{center}
\begin{tabular}{lll}
\toprule
Aspect & Object & Formula \\
\midrule
Modular (worldsheet spectrum) &
  $Z_{\mathrm{PCF}}(i)$ &
  $e^{-3\pi/2}/|\eta(i)|^6$ \\
Arithmetic (class-number formula) &
  $L(1,\chi_5)$ &
  $T/(\pi\sqrt{5})$, $\;T=2\pi\log\golden$ \\
\bottomrule
\end{tabular}
\end{center}
Both are determined by $\tau_{\mathrm{PCF}}=i$; neither can be chosen
independently.
\end{theorem}

\begin{proof}
\textit{Modular side.}
$\tau_{\mathrm{PCF}}=i$ is the unique fixed point of the S-duality
$\tau\mapsto -1/\tau$ (\cref{lem:ads-radius}).  Evaluating the Virasoro partition
function $Z(\tau) = \mathrm{Tr}\,q^{L_0-c/24}\bar q^{\bar L_0-c/24}$
at $\tau=i$ with $c=D=3$ gives $Z_{\mathrm{PCF}}(i)=e^{-3\pi/2}|\eta(i)|^{-6}$
(\cref{prop:pcf-partition}).  The modular parameter $\tau=i$ is
not integrated over (as in the free-string path integral) but is fixed
by the PCF self-duality.

\textit{Arithmetic side.}
The worldline $\Omega(\tau) = \tfrac{1}{2}e^{i\tau\log\golden}$
on $T^2_{\mathrm{PCF}}$ has the Certainty Principle
$\varepsilon(\sigma)\cdot\tau(\sigma)=\pi$ at every tower level
(the tower recurrence: $\varepsilon(\sigma) = \varepsilon_0\golden^\sigma$,
$\tau(\sigma) = \pi/(\varepsilon_0\golden^\sigma)$).  This fixes the geometric period
$T = 2\pi\log\golden$ (\cref{def:torus}).  The class-number
formula for $\mathbb{Q}(\sqrt{5})$ then gives
$L(1,\chi_5) = T/(\pi\sqrt{5})$ (\cref{thm:L1}).

\textit{Independence.}
Neither $Z_{\mathrm{PCF}}(i)$ nor $L(1,\chi_5)$ individually
determines the other: $Z_{\mathrm{PCF}}(i)$ involves $|\eta(i)|$
(a product of Gamma values at $1/4$) while $L(1,\chi_5)$ involves
$\log\golden$ (the regulator of $\mathbb{Q}(\sqrt{5})$).  These are
algebraically independent transcendental numbers.  Their common origin
is $\tau_{\mathrm{PCF}}=i$, the single geometric datum from which both
sides of $T^2_{\mathrm{PCF}}$~-- modular and arithmetic~-- are derived.
\end{proof}

% ------------------------------------------------------------------
\subsection{Bridge 3 (exact): same $\mathbb{Z}_2\times\mathbb{Z}_2$ in the Primitive Structure and here}%
\label{ssec:bridge-Z2}
% ------------------------------------------------------------------

\begin{theorem}[$\chi_5$ as a component of the Galois functor;
  Grisales Herrera~\cite{GrisalesHerrera}]%
  \label{thm:bridge3}
The character $\chi_5 = (5/\cdot)$ organising every local factor
$f_p(s)$ of $\zeta_{\mathbb{Q}(\sqrt{5})}(s)$
(\cref{prop:local-factors}) is the component $(0,1)$ of
the surjective group homomorphism
\begin{equation}\label{eq:galois-functor}
  G\colon\Ztwenty\longrightarrow\mathbb{Z}_2\times\mathbb{Z}_2,
  \qquad
  G(p\bmod 20) = \bigl(\chi_4(p),\,\chi_5(p)\bigr),
\end{equation}
where $\chi_4 = (-4/\cdot)$ is the Kronecker symbol of $\mathbb{Q}(\sqrt{-1})$.
The group $\mathbb{Z}_2\times\mathbb{Z}_2$ has exponent~$2$
(Proposition~3 of~\cite{V11}):
every element $g$ satisfies $2g = 0$, so every character is
real and self-dual.

Explicitly, the four characters of $\mathbb{Z}_2\times\mathbb{Z}_2$
classify the four quadratic fields:
\begin{center}
\begin{tabular}{cccl}
\toprule
Component & Character & Field & Split/inert criterion \\
\midrule
$(0,0)$ & trivial & $\mathbb{Q}$ & all primes split \\
$(1,0)$ & $\chi_4$ & $\mathbb{Q}(\sqrt{-1})$ & $p\equiv 1\pmod{4}$ \\
$(0,1)$ & $\chi_5$ & $\mathbb{Q}(\sqrt{5})$ & $p\equiv\pm 1\pmod{5}$ \\
$(1,1)$ & $\chi_4\chi_5$ & $\mathbb{Q}(\sqrt{-5})$ & $p\equiv 1,9\pmod{20}$ \\
\bottomrule
\end{tabular}
\end{center}
\end{theorem}

\begin{proof}
The group $\Ztwenty$ has order~$8$ with elements
$\{1,3,7,9,11,13,17,19\}$.
Define $\chi_4(a) = (-4/a)$ and $\chi_5(a) = (5/a)$ as
Kronecker symbols.  Since both are quadratic characters
(i.e.\ take values in $\{+1,-1\}$), the map
$G(a) = (\chi_4(a),\chi_5(a))$ lands in
$\mathbb{Z}_2\times\mathbb{Z}_2$.

The homomorphism property follows from the multiplicativity
of Kronecker symbols: $G(ab) = (\chi_4(ab),\chi_5(ab))
= (\chi_4(a)\chi_4(b),\chi_5(a)\chi_5(b)) = G(a)+G(b)$
(in additive notation on $\mathbb{Z}_2\times\mathbb{Z}_2$).

Surjectivity: the four fibers are
$G^{-1}(0,0)=\{1,9\}$,
$G^{-1}(1,0)=\{11,19\}$,
$G^{-1}(0,1)=\{13,17\}$,
$G^{-1}(1,1)=\{3,7\}$,
all nonempty.  (Verified by direct computation; e.g.\
$\chi_4(3)=-1$, $\chi_5(3)=-1$, so $G(3)=(1,1)$.)

Exponent~$2$: $\mathbb{Z}_2\times\mathbb{Z}_2$ has
$g+g = 0$ for all $g$, since each component is in
$\mathbb{Z}_2$.  Hence every character
$\hat\chi\colon\mathbb{Z}_2\times\mathbb{Z}_2\to\mathbb{C}^\times$
satisfies $\hat\chi^2 = 1$, i.e.\ is real and self-dual.

The correspondence between characters and quadratic
fields is classical genus theory for Dirichlet characters of
conductor dividing~$20$ (see~\cite{Neukirch}, Ch.~VII):
each primitive quadratic character $\chi_d = (d/\cdot)$
is the Artin character of $\mathbb{Q}(\sqrt{d})$, and the
four characters $\{1,\chi_4,\chi_5,\chi_4\chi_5\}$
correspond to the four fields in the table.
Verified for all primes $3\le p\le 31$ in \cref{tab:fib-split}.
\end{proof}

\begin{corollary}[Self-duality unifies this paper and the Primitive Structure]%
  \label{cor:same-Z2}
The self-duality of $\chi_5$ (which guarantees that $L(1,\chi_5)>0$
and that $\chi_5$ appears in the factorisation $\zeta = \zeta_K/L$)
follows from the exponent~$2$ of $\mathbb{Z}_2\times\mathbb{Z}_2$:
\[
  \chi_5(p)^2 = 1
  \quad\Rightarrow\quad
  \chi_5\text{ is real and self-dual.}
\]
This is the same exponent-$2$ property
that in the Primitive Structure~\cite{V11} produces the spectral constraint:
$\mathbb{Z}_2\times\mathbb{Z}_2$ having exponent~$2$ forces all characters
self-dual, which (via Hecke~\cite{Hecke}) gives $L(1-\rho,\chi)=0$ whenever
$L(\rho,\chi)=0$.

In other words: the \emph{same algebraic fact}---exponent~$2$ of
$\mathbb{Z}_2\times\mathbb{Z}_2$---simultaneously produces
the spectral constraint of the Primitive Structure~\cite{V11} and makes $\chi_5$ real and self-dual
so that $L(1,\chi_5) = 2\log\golden/\sqrt{5} > 0$ is the class-number
residue (this paper).
\end{corollary}

The categorical reading of this bridge is the span of
\cref{def:pcf-span}: $G$ is the common source of
both the arithmetic branch (which produces $\zeta(2k+1)$ via
the Euler product colimit) and the spectral branch (which
produces $\mu_3=1/2$ via $|\ker G|=2$).
The exponent-$2$ property operates in both branches
simultaneously: in the arithmetic branch it makes $\chi_5$
self-dual, ensuring $L(1,\chi_5)>0$; in the spectral branch
it forces all characters self-dual, enabling Hecke's companion
bound.

% ------------------------------------------------------------------
\subsection{Summary: what the three bridges prove}%
\label{ssec:bridges-summary}
% ------------------------------------------------------------------

\begin{enumerate}
\item \textbf{Bridge~1 (exact, self-contained):}
  $S_{\mathrm{BH}} = k_B\lambda_{\log}\tfrac{\sqrt{5}}{2}L(1,\chi_5)$.
  Derived from $|\Omega|=1/2$ (\cref{prop:worldline-entropy})
  and the Mersenne bridge; error $= 0$ at 50 digits.
\item \textbf{Bridge~2 (structural):}
  $Z_{\mathrm{PCF}}(i)$ and $L(1,\chi_5)$ both arise from $\tau_{\mathrm{PCF}}=i$,
  through modular and arithmetic aspects of $T^2_{\mathrm{PCF}}$ respectively.
\item \textbf{Bridge~3 (exact, self-contained):}
  $\chi_5$ is the $(0,1)$-component of $G\colon\Ztwenty\to\mathbb{Z}_2\times\mathbb{Z}_2$
  (\cref{thm:bridge3});
  exponent~$2$ makes $\chi_5$ self-dual here and produces the spectral
  constraint of the Primitive Structure~\cite{V11}.
\end{enumerate}

Together, the three bridges formalise the open problem of the Primitive Structure~\cite{V11}:
the structural content of $\mu_3^2=1/4$ (equivalently $1/(4G_N)=\mu_3=1/2$) is the class-number
formula $L(1,\chi_5)=2\log\golden/\sqrt{5}$ expressed as
$S_{\mathrm{BH}}/k_B = \lambda_{\log}(\sqrt5/2)L(1,\chi_5)$
(Bridge~1); the Virasoro central charge $c=3$ arises from the same
$S_3$ symmetry of $T^2_{\mathrm{PCF}}$ that fixes $\tau_{\mathrm{PCF}}=i$
and hence both $Z_{\mathrm{PCF}}(i)$ and $L(1,\chi_5)$ (Bridge~2);
and the $\chi_5$-character organising odd zeta values is the
$(0,1)$-component of the Galois functor whose exponent-$2$ property
produces both the spectral constraint of the Primitive Structure~\cite{V11} and the
self-duality of $L(1,\chi_5)$ (Bridge~3).

% ------------------------------------------------------------------
\subsection{The unifying generator}\label{ssec:unifying}
% ------------------------------------------------------------------

We collect the results of \cref{sec:background,ssec:bridge-Z2}
into a single statement.  The relation $\golden^2 = \golden + 1$
generates every structure in this paper
(\cref{thm:unifying} below), and the dimensional chain from
1D through 2D to 3D follows as a consequence of the Euler product
over the primes (\cref{cor:S3-from-euler}).

\begin{theorem}[$\golden^2 = \golden+1$ generates the complete structure]%
  \label{thm:unifying}
The single relation $\golden^2 = \golden + 1$ generates the complete
arithmetic-geometric structure of $T^2_{\mathrm{PCF}}$ and determines
every object in this paper:
\begin{enumerate}
\item[\textup{(1)}] \textbf{Frobenius:} $\golden^p = F_p\golden + F_{p-1}$
  for every prime $p$, with splitting type
  $\chi_5(p)\equiv F_p\pmod{p}$
  (\cref{prop:fib-split});
\item[\textup{(2)}] \textbf{Mersenne:} $\golden^{\lambda_{\log}} = 2$,
  $\lambda_{\log} = \log 2/\log\golden$
  (\cref{thm:mersenne});
\item[\textup{(3)}] \textbf{Pentagon:} $\golden = 2\cos(\pi/5)$,
  connecting to $\mathbb{Q}(\zeta_5)$
  (\cref{prop:pentagon});
\item[\textup{(4)}] \textbf{Galois functor:}
  $G\colon\Ztwenty\to\mathbb{Z}_2\times\mathbb{Z}_2$ with
  exponent~$2$, $|\ker G| = 2$, and four self-dual characters
  (\cref{thm:bridge3});
\item[\textup{(5)}] \textbf{Torus:}
  $\Omega(\tau) = \tfrac{1}{2}e^{i\tau\log\golden}$ closes
  with period $T = 2\pi\log\golden$
  (\cref{def:torus});
\item[\textup{(6)}] \textbf{Norm:} $|\Omega| = 1/2$
  (\cref{prop:worldline-entropy});
\item[\textup{(7)}] \textbf{Euler product:}
  $\zeta_{\mathbb{Q}(\sqrt{5})}(s) = \prod_p f_p(s)$ and
  $\zeta(s) = \zeta_{\mathbb{Q}(\sqrt{5})}(s)/L(s,\chi_5)$
  (\cref{thm:factorisation});
\item[\textup{(8)}] \textbf{Odd zeta values:} $\zeta(2k+1)$
  structurally determined
  (\cref{thm:pentagon-odd});
\item[\textup{(9)}] \textbf{Categorical framework:} the Frobenius
  functor $\Psi$, the Euler product colimit, the Dedekind natural
  isomorphism, and the arithmetic--spectral span are all generated
  by $\golden^2=\golden+1$ through the constructions of
  \cref{sec:categorical-framework}.
\end{enumerate}
The only inputs beyond $\golden^2 = \golden + 1$ are classical
number theory (\cite{Neukirch}, \cite{Washington}) and the
complexification $e^{ix}$ (standard analysis).
\end{theorem}

\begin{proof}
Every item derives from $\golden^2 = \golden + 1$ by an explicit
chain already established in this paper:
(1)~follows by induction: $\golden^{n+2} = \golden^{n+1}+\golden^n$
reproduces the Fibonacci recurrence in the coefficients;
(2)~follows from $e^{(\log 2/\log\golden)\log\golden} = e^{\log 2} = 2$;
(3)~follows from $\cos(\pi/5) = \golden/2$ (pentagonal geometry);
(4)~follows from (1) and (3): $\chi_5 = (5/\cdot)$ is the Artin
character of $\mathbb{Q}(\sqrt{5})$, the field generated by $\golden$,
and $\chi_4 = (-4/\cdot)$ classifies primes mod~$4$; together they
define $G$ on $\Ztwenty$;
(5)~follows from complexifying $\golden^\sigma$ via
$\tau\mapsto e^{i\tau\log\golden}$;
(6)~follows from $|e^{ix}|=1$;
(7)~follows from the Dedekind factorisation of $\mathbb{Q}(\sqrt{5})$,
whose discriminant $\Delta = 5$ and class number $h = 1$ are
determined by $\golden^2 = \golden + 1$;
(8)~follows from (1)--(7) combined;
(9)~follows from (1)--(8) by the constructions of
\cref{sec:categorical-framework}: the functoriality of $\Psi$
is~(1), the colimit is~(7), the natural isomorphism is
(7)$+$(4), and the span is their conjunction.
\end{proof}

\begin{lemma}[AdS radius from S-duality]\label{lem:ads-radius}
The modular parameter $\tau_{\mathrm{PCF}} = i$
(\cref{def:torus}) is the unique fixed point of
S-duality $S\colon\tau\mapsto -1/\tau$ (since $-1/i = i$).
The AdS radius $\ell$ transforms as $\ell\mapsto 1/\ell$ under~$S$.
At the fixed point: $\ell = 1/\ell$, giving $\ell = 1$.
This is not a convention but the unique value compatible with
the S-autoduality inherited from
$\golden\leftrightarrow -1/\golden$
(\cref{prop:phi-unique}).
\end{lemma}

\begin{proof}
$S\colon\tau\mapsto -1/\tau$ acts on $\tau_{\mathrm{PCF}} = i$ by
$-1/i = i$, so $\tau_{\mathrm{PCF}}$ is the unique fixed point of~$S$
on the upper half-plane.
Under the standard identification $\tau = i\ell/\ell_s$ with
string length $\ell_s$, the S-map exchanges $\ell\leftrightarrow 1/\ell$
(in units $\ell_s = 1$).
At the fixed point: $\ell = 1/\ell$, so $\ell^2 = 1$; since
$\ell > 0$, we get $\ell = 1$.
The Galois conjugation $\golden\leftrightarrow\bar\golden = -1/\golden$
is the arithmetic source: $N(\golden) = \golden\bar\golden = -1$
is the norm that generates the S-duality at the level of~$\Rpcf$.
\end{proof}

\begin{corollary}[$S_3$, $E^3$, and $c = 3$ from the Euler product]%
  \label{cor:S3-from-euler}
The Euler product $\zeta_{\mathbb{Q}(\sqrt{5})}(s) = \prod_p f_p(s)$ is the collective
iteration of the Frobenius self-measurement
$\golden^p = F_p\golden+F_{p-1}$ over all primes, converging
on~$\mathbb{C}$.  The Galois functor
$G\colon\Ztwenty\to\mathbb{Z}_2\times\mathbb{Z}_2$
(\cref{thm:bridge3}) has $|\ker G| = 2$, producing a
contraction factor $\mu_2/\mu_3 = 2$ from the binary level
($\mu_2 = 1$) to the ternary level ($\mu_3 = 1/2 = |\Omega|$).

The spectral constraints (the ternary fragment $n=3$ is the unique
arity where the Lawvere~\cite{Lawvere} obstruction is simultaneously effective
$\mu_n < 1$ and non-degenerate $\mu_n > 0$; the Moran equation
$n\mu^d = 1$ with $n = 3$, $d = 1$ gives $\mu = 1/3$, but the
correct counting uses the spectral product $\sigma\mu$):
\begin{equation}\label{eq:spectral-constraints}
  \sigma + \mu = 2, \qquad \sigma\,\mu = \tfrac{3}{4}
\end{equation}
force $\mu = 1/2$, $\sigma = 3/2$ uniquely (the quadratic
$\sigma^2 - 2\sigma + 3/4 = 0$ has roots $3/2$ and $1/2$;
only $\mu = 1/2 < 1$ gives an effective obstruction).

The tripartite norm decomposition
(the three eigenvalues
$\omega^k/2$ form an equilateral triangle inscribed in $|z|=1/2$;
$|C| = 1$ is the identity component at the real eigenvalue
$\lambda_0 = 1/2$; $|F| = |1-\omega|/2 = \sqrt{3}/2$ is the
side length; $|P| = 1/\sqrt{3}$ is the inradius-to-side ratio):
\begin{equation}\label{eq:tripartite-norm}
  \|P\| \cdot \|C\| \cdot \|F\|
  = \frac{1}{\sqrt{3}}\cdot 1 \cdot\frac{\sqrt{3}}{2}
  = \frac{1}{2} = |\Omega|
\end{equation}
exhibits $S_3$ symmetry: the three eigenvalues
$\omega^k/2$ ($\omega = e^{2\pi i/3}$, $k = 0,1,2$) form an
equilateral triangle of radius~$1/2$ in~$\mathbb{C}$, permuted
by the symmetric group $S_3$.  This $S_3$ determines:
the golden coupling $z = \golden y$
(\cref{prop:dim-chain}),
the space $E^3 = \{(x,y,\golden y)\}$, and the central charge
$c = 3\ell/(2G_N)|_{\ell=1,\,G_N=1/2} = 3$
(\cref{lem:ads-radius}).

The chain is therefore:
\[
  \underbrace{\golden^2 = \golden+1}_{\textup{generator}}
  \;\to\;
  \underbrace{\prod_p f_p(s)}_{\textup{Euler product on }\mathbb{C}}
  \;\to\;
  \underbrace{|\ker G|=2}_{\textup{contraction}}
  \;\to\;
  \underbrace{\mu_3 = \tfrac{1}{2}}_{\textup{norm}}
  \;\to\;
  \underbrace{S_3}_{\textup{symmetry}}
  \;\to\;
  \underbrace{E^3,\; c=3.}_{\textup{geometry}}
\]
The third dimension and its properties are not additional
data---they are consequences of the Euler product of
$\golden^2 = \golden + 1$ over the primes.
Both are verified in Lean~$4$ with $0$~\texttt{sorry}:
\texttt{norm\_Omega\_is\_half}, \texttt{spectral\_uniqueness},
and \texttt{galois\_functor\_complete} (\cite{V11}).
\end{corollary}

\begin{proof}
$|\ker G| = 2$: \cref{thm:bridge3} gives
$G^{-1}(0,0) = \{1,9\}$.
Contraction: $\mu_2/\mu_3 = 1/(1/2) = 2 = |\ker G|$.
Spectral uniqueness: substituting $\mu = 2-\sigma$ into
$\sigma\mu = 3/4$ gives $\sigma^2 - 2\sigma + 3/4 = 0$;
roots $\sigma = (2\pm 1)/2 \in \{3/2,\,1/2\}$.
Only $\sigma = 3/2$ gives $\mu = 1/2 < 1$.
Tripartite norm: $(1/\sqrt{3})\cdot 1\cdot(\sqrt{3}/2) = 1/2$.
$S_3$ symmetry: $\omega = e^{2\pi i/3}$ satisfies $\omega^3 = 1$;
the three eigenvalues $\{\omega^0/2,\,\omega^1/2,\,\omega^2/2\}$
are permuted by the symmetric group on three elements.
Central charge: The holographic identification $G_N = \mu_3$
follows from $S = A/(4G_N)$ with $\mu_3^2 = 1/4$: the Planck
area factor $1/4$ is the square of
the spectral parameter derived from $S_3$ symmetry, giving
$G_N = \mu_3 = 1/2$.
Then $c = 3\ell/(2G_N) = 3\cdot 1/(2\cdot 1/2) = 3$
(Brown--Henneaux with $\ell = 1$ (\cref{lem:ads-radius})).
\end{proof}

\begin{lemma}[Spectral real parts]\label{lem:spectral-re}
The three eigenvalues $\omega^k/2$ ($\omega = e^{2\pi i/3}$,
$k = 0,1,2$) have real parts:
\[
  \mathrm{Re}(\omega^0/2) = \tfrac{1}{2}, \qquad
  \mathrm{Re}(\omega^1/2) = \mathrm{Re}(\omega^2/2) = -\tfrac{1}{4}.
\]
The unique eigenvalue with $\mathrm{Re} > 0$ has
$\mathrm{Re} = 1/2 = \mu_3$.
\end{lemma}

\begin{proof}
$\omega^0 = 1$, so $\mathrm{Re}(\omega^0/2) = 1/2$.
$\omega = e^{2\pi i/3} = -1/2 + i\sqrt{3}/2$, so
$\mathrm{Re}(\omega/2) = -1/4$.
$\omega^2 = \bar\omega$, so $\mathrm{Re}(\omega^2/2) = -1/4$.
Only $k = 0$ gives $\mathrm{Re} > 0$, and
$\mathrm{Re}(\omega^0/2) = 1/2 = \mu_3$
(\cref{prop:worldline-entropy}).
Verified in Lean~$4$: \texttt{eigenvalue\_half} ($0$~\texttt{sorry}).
\end{proof}

\begin{proposition}[Spectral structure of the Euler product]%
  \label{prop:spectral-id}
The map $G\colon\Ztwenty\to\mathbb{Z}_2\times\mathbb{Z}_2$
that organises the Euler product of $\zeta(s)$ also determines
the eigenvalue structure of $\OmH$:
\begin{enumerate}
\item[\textup{(1)}] $G$ classifies every prime $p > 5$ by
  splitting type in $\mathbb{Z}[\golden]$
  (\cref{thm:bridge3},
  \cref{prop:fib-split}).
\item[\textup{(2)}] The splitting type determines the local
  Euler factor $f_p(s)$
  (\cref{prop:local-factors}).
\item[\textup{(3)}] The Euler product
  $\zeta_{\mathbb{Q}(\sqrt{5})}(s) = \prod_p f_p(s)$ and the factorisation
  $\zeta(s) = \zeta_{\mathbb{Q}(\sqrt{5})}(s)/L(s,\chi_5)$ hold for
  $\mathrm{Re}(s)>1$ (\cref{thm:factorisation};
  Euler~\cite{Euler}).
\item[\textup{(4)}] Analytic continuation determines $\zeta(s)$
  on all of $\mathbb{C}$ from its values on $\mathrm{Re}(s)>1$
  (identity theorem for Dirichlet series; Riemann~\cite{Riemann};
  \cref{rmk:euler-determines}).
\item[\textup{(5)}] The same $G$ determines
  $|\ker G|=2\to\mu_3=1/2\to S_3\to$
  eigenvalues $\omega^k/2$
  (\cref{cor:S3-from-euler}); the unique
  eigenvalue with $\mathrm{Re}>0$ has $\mathrm{Re}=1/2$
  (\cref{lem:spectral-re}).
\end{enumerate}
Steps (1)--(4) and (5) are two readings of the same map $G$:
one produces $\zeta(s)$, the other produces the spectral
data of $\OmH$.
\end{proposition}

\begin{proof}
Steps (1)--(3) are
\cref{thm:bridge3},
\cref{prop:local-factors}, and
\cref{thm:factorisation} respectively.
Step (4) is the identity theorem for Dirichlet series
(\cref{rmk:euler-determines}).
Step (5) is \cref{cor:S3-from-euler} and
\cref{lem:spectral-re}.
\end{proof}

\begin{remark}[Discharge of the bridge axiom]\label{rmk:bridge-resolution}
\Cref{prop:spectral-id} makes explicit the mechanism that
the Primitive Structure~\cite{V11} formalises as \texttt{bridge\_axiom}
in Lean~4.  In the Primitive Structure, Step~4 of the Spectral Embedding
(Proposition~5 of~\cite{V11}) states that zeros of $L$-functions factorising
through~$\Ztwenty$ inherit the spectral constraint
$\mathrm{Re}(\rho)\in\mathrm{Re}(\mathrm{spec}(\OmH))
= \{1/2,\,-1/4\}$.  The Lean formalisation encodes this as an axiom
(\texttt{bounded\_by\_mu}) because Mathlib currently lacks Dirichlet
series and the spectral mapping theorem for Banach algebras.

The present paper provides the mathematical justification:
\begin{enumerate}
\item $G$ classifies every prime $p>5$ by splitting type in
  $\mathbb{Z}[\golden]$, determining every local Euler factor $f_p(s)$
  (\cref{prop:local-factors}).
\item The Euler product $\prod_p f_p(s)$ determines $\zeta(s)$ on
  $\mathrm{Re}(s)>1$ (Euler~\cite{Euler}), and hence on all
  of~$\mathbb{C}$ by analytic continuation (Riemann~\cite{Riemann}).
\item The \emph{same} $G$ determines $|\ker G|=2\to\mu_3=1/2\to S_3\to$
  eigenvalues $\omega^k/2$ (\cref{cor:S3-from-euler}).
\end{enumerate}
The bridge is therefore not an additional hypothesis: it is the conjunction
of steps (1)--(2) (classical number theory) with step~(3) (verified in
Lean).  The \texttt{bridge\_axiom} of the Primitive Structure is the
formalisation of the fact that the same map~$G$ simultaneously produces
the analytic object ($\zeta$, via the Euler product) and the spectral
constraint ($\mu_3=1/2$, via the kernel).
\end{remark}

\begin{remark}[Scope]\label{rmk:scope}
This paper establishes the isomorphism between
$\golden^2=\golden+1$ and the components of~$\zeta(s)$ at every level
of the Euler product (\cref{prop:spectral-id}), thereby
discharging the \texttt{bridge\_axiom}
(\cref{rmk:bridge-resolution}).

What is \emph{not} addressed here is the \emph{functorial transport}
of the eigenvalue constraint $\mathrm{Re}>0\Rightarrow\mathrm{Re}=1/2$
from $\OmH$ to~$\zeta$ across the full critical strip---this requires the
forgetful functor chain
$C^*$-$\mathbf{Alg}\to\mathbf{Hilb}\to\mathbf{Met}\to\mathbf{Set}$
and the categorical limit theorem, which are established in the Primitive
Structure~\cite{V11} (Lean: \texttt{categorical\_limit\_theorem},
\texttt{Omega\_hat\_unique\_positive\_re}, $0$~\texttt{sorry}).
\end{remark}

\begin{remark}[Classical inputs]\label{rmk:classical-inputs}
The chain of \cref{prop:spectral-id} uses,
beyond the results established in this paper, three classical theorems:
\begin{enumerate}
\item \textbf{Euler product identity} (Euler~\cite{Euler}):
  $\prod_p(1-p^{-s})^{-1} = \sum_{n=1}^{\infty}n^{-s}$ for
  $\mathrm{Re}(s)>1$.  Proof: unique factorisation of $\mathbb{N}$.
\item \textbf{Identity theorem for Dirichlet series}
  (Riemann~\cite{Riemann}):
  the coefficients of a Dirichlet series determine the function
  on all of $\mathbb{C}$ via analytic continuation.
  Therefore $G$ determines $\zeta$ everywhere.
\item \textbf{Isomorphism transport}:
  if $A\cong B$ and $P(A)$, then $P(B)$.
  The identification runs through five levels
  (\cref{prop:spectral-id}): primes, splitting types,
  Euler factors, $\zeta$, and simultaneously
  $|\ker G|\to\mu_3\to S_3\to$ eigenvalues.
\end{enumerate}
All three are standard results with known proofs; none is specific
to the PCF framework.  Their conjunction is what the Primitive
Structure~\cite{V11} formalises as \texttt{bridge\_axiom}
(\cref{rmk:bridge-resolution}).
\end{remark}

%%% ====================================================================
